
\subsection{Identification}

For a beam formulation combining both problems, the complete rigid
body freedom can be parametrized as:

\[
\hat{\bm{u}} = \underbrace{\bm{\beta}_u + c_\vartheta(\FrameBasis \times \bm{r}) + (\FrameBasis \times \bm{\beta}_\kappa) \times \bm{r}}_{\text{rigid at } x=0} + x\underbrace{[a_\varepsilon \FrameBasis + \bm{\beta}_\gamma + (a_\vartheta \FrameBasis - \FrameBasis \times \bm{a}_\Omega) \times \bm{r}]}_{\text{rate terms}} + \text{(higher order terms)}
\tag{V1}
\]

Your choice of reference curve determines the relationship:
\[
\begin{aligned}
\bm{u}(\reflenx) &= \bm{\beta}_u + x(a_\varepsilon \FrameBasis + \bm{\beta}_\gamma) + \text{(higher order)} \\
\bm{\theta}(\reflenx) &= \bm{\beta}_\theta + \reflenx\, (a_\vartheta \FrameBasis - \FrameBasis \times \bm{a}_\Omega) + \text{(higher order)}
\end{aligned}
\tag{V1}
\]

(V3) If we write the general solution as:
\[
\hat{\bm{u}} = \bm{u}(\reflenx) + \bm{\theta}(\reflenx) \times \bm{r} + \text{(warping)}
\tag{V3}
\]

Then choosing \(\bm{u}(\reflenx)\) to include
\(x\bm{\beta}_\gamma\) and \(\bm{\theta}(\reflenx)\) to include
\(-\FrameBasis \times \bm{\beta}_\gamma\) 
produces the constant shear strain
\[
\bm{\gamma} = \bm{u}' + \FrameBasis \times \bm{\theta} = \bm{\beta}_\gamma
\tag{V3}
\]
Set:
\begin{equation}
\begin{aligned}
\bm{u}(\reflenx) &= \bm{\beta}_u + x(a_\varepsilon \FrameBasis + \bm{\beta}_\gamma) - \tfrac{1}{2}x^2 \bm{a}_\Omega - \tfrac{1}{6}x^3 \bm{b}_{\Section} - \tfrac{1}{2}x^2(\FrameBasis \otimes \bar{\bm{r}})\bm{b}_{\Section} \\
\bm{\theta}(\reflenx) &= \bm{\beta}_\theta + x(a_\vartheta + c_\vartheta)\FrameBasis - x(\FrameBasis \times \bm{a}_\Omega) - \tfrac{1}{2}x^2(\FrameBasis \times \bm{b}_{\Section}) - \FrameBasis \times \bm{\beta}_\gamma
\end{aligned}
\tag{V2/3}
\end{equation}


where \(\bm{\beta}_\gamma\) is determined by a suitable fitting
criterion (energy conjugacy), and this determines how the transverse
components of \(\bm{\beta}_\theta\) are split between the
reference curve displacement and rotation.
%
Note that \(\bm{\beta}_\gamma\) appears in both
\(\bm{u}'\) and \(\FrameBasis \times \bm{\theta}\), producing the
constant shear strain
\(\bm{\gamma} = \bm{\beta}_\gamma\) that you use as your
fitting parameter to relate
\(\bm{b}_{\Section} = \bm{\Xi} \bm{\gamma}\).

\subsection{Key Insights: (V1)}

\begin{enumerate}
\def\labelenumi{\arabic{enumi}.}
\item
  \textbf{The \(\reflenx\)-independence} in Problem I refers to fields that
  don't vary with the axial coordinate within each term's coefficient.
  This carries to Problem II for the coefficients of each power of
  \(\reflenx\).
\end{enumerate}

